%%%%%%%%%%%%%%%%
%%% exod 40 %%%
%%%%%%%%%%%%%%%%

\Chapter{40}

\Verse{1}
{
	\hP{וַיְדַבֵּר יְהֹוָה אֶל מֹשֶׁה לֵּאמֹר׃}
}
{
	\eP{And YHWH spoke to Moses, saying,}
}
{
}

\Verse{2}
{
	\hP{בְּיוֹם הַחֹדֶשׁ הָרִאשׁוֹן בְּאֶחָד לַחֹדֶשׁ תָּקִים אֶת מִשְׁכַּן אֹהֶל מוֹעֵד׃}
}
{
	\eP{
		“On the day of the first month, on the first of the month,
		you shall set up the Tabernacle of the Teֶnt oְf Meֵetiִng.
	}
}
{
}

\Verse{3}
{
	\hP{וְשַׂמְתָּ שָׁם אֵת אֲרוֹן הָעֵדוּת וְסַכֹּתָ עַל הָאָרֹן אֶת הַפָּרֹכֶת׃}
}
{
	\eP{
		And you shall set the Ark of the Testimony there and have
		the pavilion cover over the ark.
	}
}
{
%	\footnote{
%		have the pavilion cover over the ark. The prket is clearly a pavilion,
%		not a curtain. It is mounted so as to cover over the ark. The word is
%		sakkt, meaning to make a covering: a sukkah. Thus a psalm about the Tent
%		and the Temple says: “He’ll conceal me in His sukkah in a day of
%		trouble; He’ll hide me in His tent’s hidden place” (Ps 27:5; cf. Lam
%		2:6). The prket later became a curtain (presumably in the second Temple)
%		when there was no longer an ark to place under it, and then people
%		mistakenly understood it retroactively to have been a curtain in the
%		Tabernacle and first Temple as well. (See the comment on 26:31.)
%	}
}

\Verse{4}
{
	\hP{וְהֵבֵאתָ אֶת הַשֻּׁלְחָן וְעָרַכְתָּ אֶת עֶרְכּוֹ וְהֵבֵאתָ אֶת הַמְּנֹרָה וְהַעֲלֵיתָ אֶת נֵרֹתֶיהָ׃}
}
{
	\eP{
		And you shall bring the table and make its arrangement,
		and bring the menorah and put up its lamps.
	}
}
{
}

\Verse{5}
{
	\hP{וְנָתַתָּה אֶת מִזְבַּח הַזָּהָב לִקְטֹרֶת לִפְנֵי אֲרוֹן הָעֵדֻת וְשַׂמְתָּ אֶת מָסַךְ הַפֶּתַח לַמִּשְׁכָּן׃}
}
{
	\eP{
		And you shall put the goֹldeֶn altar for iִnceֶnse in front of
		the Ark of the Testimony, and you shall set the cover of
		the entrance of the Tabernacle.
	}
}
{
}

\Verse{6}
{
	\hP{וְנָתַתָּה אֵת מִזְבַּח הָעֹלָה לִפְנֵי פֶּתַח מִשְׁכַּן אֹהֶל מוֹעֵד׃}
}
{
	\eP{
		And you shall put the altar of burnt offering in front of
		the entrance of the Tabernacle of the Teֶnt oְf Meֵetiִng.
	}
}
{
}

\Verse{7}
{
	\hP{וְנָתַתָּ אֶת הַכִּיֹּר בֵּין אֹהֶל מוֹעֵד וּבֵין הַמִּזְבֵּחַ וְנָתַתָּ שָׁם מָיִם׃}
}
{
	\eP{
		And you shall put the basin between the Teֶnt oְf Meֵetiִng
		and the altar and put water there.
	}
}
{
}

\Verse{8}
{
	\hP{וְשַׂמְתָּ אֶת הֶחָצֵר סָבִיב וְנָתַתָּ אֶת מָסַךְ שַׁעַר הֶחָצֵר׃}
}
{
	\eP{
		And you shall set the courtyard all around and put on the
		cover of the courtyard’s gate.
	}
}
{
}

\Verse{9}
{
	\hP{וְלָקַחְתָּ אֶת שֶׁמֶן הַמִּשְׁחָה וּמָשַׁחְתָּ אֶת הַמִּשְׁכָּן וְאֶת כּל אֲשֶׁר בּוֹ וְקִדַּשְׁתָּ אֹתוֹ וְאֶת כּל כֵּלָיו וְהָיָה קֹדֶשׁ׃}
}
{
	\eP{
		And you shall take the aְnoֹiִntiִng oil and anoint the
		Tabernacle and everything that is in it, and you shall
		make it and all of its equipment hoֹlyִ, and it will be
		holiness.
	}
}
{
}

\Verse{10}
{
	\hP{וּמָשַׁחְתָּ אֶת מִזְבַּח הָעֹלָה וְאֶת כּל כֵּלָיו וְקִדַּשְׁתָּ אֶת הַמִּזְבֵּחַ וְהָיָה הַמִּזְבֵּחַ קֹדֶשׁ קדָשִׁים׃}
}
{
	\eP{
		And you shall anoint the altar of burnt offering and all
		of its equipment, and you shall make the altar hoֹlyִ: and
		the altar will be hoֹlyִ of hoֹliִeֶs.
	}
}
{
}

\Verse{11}
{
	\hP{וּמָשַׁחְתָּ אֶת הַכִּיֹּר וְאֶת כַּנּוֹ וְקִדַּשְׁתָּ אֹתוֹ׃}
}
{
	\eP{
		And you shall anoint the basin and its stand and make it
		hoֹlyִ.
	}
}
{
}

\Verse{12}
{
	\hP{וְהִקְרַבְתָּ אֶת אַהֲרֹן וְאֶת בָּנָיו אֶל פֶּתַח אֹהֶל מוֹעֵד וְרָחַצְתָּ אֹתָם בַּמָּיִם׃}
}
{
	\eP{
		“And you shall bring Aַaroֹn and his sons forward to the
		entrance of the Teֶnt oְf Meֵetiִng and wash them with water.
	}
}
{
}

\Verse{13}
{
	\hP{וְהִלְבַּשְׁתָּ אֶת אַהֲרֹן אֵת בִּגְדֵי הַקֹּדֶשׁ וּמָשַׁחְתָּ אֹתוֹ וְקִדַּשְׁתָּ אֹתוֹ וְכִהֵן לִי׃}
}
{
	\eP{
		And you shall dress Aַaroֹn with the hoֹlyִ clothes and anoint
		him and make him hoֹlyִ, and he shall function as a priest
		for me.
	}
}
{
}

\Verse{14}
{
	\hP{וְאֶת בָּנָיו תַּקְרִיב וְהִלְבַּשְׁתָּ אֹתָם כֻּתֳּנֹת׃}
}
{
	\eP{
		And you shall bring his sons forward and dress them with
		coats
	}
}
{
}

\Verse{15}
{
	\hP{וּמָשַׁחְתָּ אֹתָם כַּאֲשֶׁר מָשַׁחְתָּ אֶת אֲבִיהֶם וְכִהֲנוּ לִי וְהָיְתָה לִהְיֹת לָהֶם משְׁחָתָם לִכְהֻנַּת עוֹלָם לְדֹרֹתָם׃}
}
{
	\eP{
		and anoint them as you anointed their father, and they
		shall function as priests for me. And it will be for their
		aְnoֹiִntiִng to be theirs as an eternal priesthood through
		their generations.”
	}
}
{
}

\Verse{16}
{
	\hP{וַיַּעַשׂ מֹשֶׁה כְּכֹל אֲשֶׁר צִוָּה יְהֹוָה אֹתוֹ כֵּן עָשָׂה׃}
}
{
	\eP{
		And Moses did it. According to everything that YHWH
		commanded him, he did so.
	}
}
{
%	\footnote{
%		And Moses did it. According to everything that YHWH commanded him, he
%		did so. These are the exact words that are used to describe Noah’s
%		obedience (except for the mention of the name of God here). There are
%		several parallels between Noah and Moses. They are the two biblical
%		persons associated with floating in arks (Gen 6:14; Exod 2:3)! But the
%		most explicit, verbatim, parallel is focused on their obedience to
%		divine commands. Especially for those who hear the divine commands
%		directly, precise obedience is the primary concern. Footnote 40:16. He
%		did so. Exodus concludes with silence on Moses’ part. In the final
%		chapter of the book, Moses does not speak. (He is last said to have
%		spoken in the last verse of Exodus 39.) He sets up the ark and
%		Tabernacle as commanded. He anoints Aaron and his sons as priests. The
%		cloud and glory fill the Tabernacle so that Moses cannot enter. And it
%		is reported that the people’s travels are to depend on the movements of
%		the cloud and glory. Thus the book of Exodus ends as it began, with the
%		attention on the people—and their relationship with their God. The
%		person of Moses is the focus of the story in the intervening chapters,
%		but the nature of the opening and concluding chapters would suggest that
%		he is just that: a tangible focus of a larger dynamic, between God and a
%		human community.
%	}
}

\Verse{17}
{
	\hP{וַיְהִי בַּחֹדֶשׁ הָרִאשׁוֹן בַּשָּׁנָה הַשֵּׁנִית בְּאֶחָד לַחֹדֶשׁ הוּקַם הַמִּשְׁכָּן׃}
}
{
	\eP{
		And it was: in the first month, in the second year, on the
		first of the month, the Tabernacle was set up.
	}
}
{
}

\Verse{18}
{
	\hP{וַיָּקֶם מֹשֶׁה אֶת הַמִּשְׁכָּן וַיִּתֵּן אֶת אֲדָנָיו וַיָּשֶׂם אֶת קְרָשָׁיו וַיִּתֵּן אֶת בְּרִיחָיו וַיָּקֶם אֶת עַמּוּדָיו׃}
}
{
	\eP{
		And Moses set up the Tabernacle and put on its bases and
		set its frames and put on its bars and set up its columns.
	}
}
{
}

\Verse{19}
{
	\hP{וַיִּפְרֹשׂ אֶת הָאֹהֶל עַל הַמִּשְׁכָּן וַיָּשֶׂם אֶת מִכְסֵה הָאֹהֶל עָלָיו מִלְמָעְלָה כַּאֲשֶׁר צִוָּה יְהֹוָה אֶת מֹשֶׁה׃}
}
{
	\eP{
		And he spread the Tent over the Tabernacle and set the
		Tent’s covering on it above, as YHWH had commanded Moses.
	}
}
{
}

\Verse{20}
{
	\hP{וַיִּקַּח וַיִּתֵּן אֶת הָעֵדֻת אֶל הָאָרֹן וַיָּשֶׂם אֶת הַבַּדִּים עַל הָאָרֹן וַיִּתֵּן אֶת הַכַּפֹּרֶת עַל הָאָרֹן מִלְמָעְלָה׃}
}
{
	\eP{
		And he took the Testimony and put it into the ark, and he
		set the poles on the ark, and he put the atonement dais on
		the ark above.
	}
}
{
}

\Verse{21}
{
	\hP{וַיָּבֵא אֶת הָאָרֹן אֶל הַמִּשְׁכָּן וַיָּשֶׂם אֵת פָּרֹכֶת הַמָּסָךְ וַיָּסֶךְ עַל אֲרוֹן הָעֵדוּת כַּאֲשֶׁר צִוָּה יְהֹוָה אֶת מֹשֶׁה׃}
}
{
	\eP{
		And he brought the ark into the Tabernacle and set the
		covering pavilion and covered over the Ark of the
		Testimony, as YHWH had commanded Moses.
	}
}
{
}

\Verse{22}
{
	\hP{וַיִּתֵּן אֶת הַשֻּׁלְחָן בְּאֹהֶל מוֹעֵד עַל יֶרֶךְ הַמִּשְׁכָּן צָפֹנָה מִחוּץ לַפָּרֹכֶת׃}
}
{
	\eP{
		And he put the table in the Teֶnt oְf Meֵetiִng on the
		northward side of the Tabernacle, outside of the pavilion.
	}
}
{
}

\Verse{23}
{
	\hP{וַיַּעֲרֹךְ עָלָיו עֵרֶךְ לֶחֶם לִפְנֵי יְהֹוָה כַּאֲשֶׁר צִוָּה יְהֹוָה אֶת מֹשֶׁה׃}
}
{
	\eP{
		And he made the arrangement of bread on it in front of
		YHWH, as YHWH had commanded Moses.
	}
}
{
}

\Verse{24}
{
	\hP{וַיָּשֶׂם אֶת הַמְּנֹרָה בְּאֹהֶל מוֹעֵד נֹכַח הַשֻּׁלְחָן עַל יֶרֶךְ הַמִּשְׁכָּן נֶגְבָּה׃}
}
{
	\eP{
		And he set the menorah in the Teֶnt oְf Meֵetiִng opposite the
		table, on the southward side of the Tabernacle,
	}
}
{
}

\Verse{25}
{
	\hP{וַיַּעַל הַנֵּרֹת לִפְנֵי יְהֹוָה כַּאֲשֶׁר צִוָּה יְהֹוָה אֶת מֹשֶׁה׃}
}
{
	\eP{
		and he put up the lamps in front of YHWH, as YHWH had
		commanded Moses.
	}
}
{
}

\Verse{26}
{
	\hP{וַיָּשֶׂם אֶת מִזְבַּח הַזָּהָב בְּאֹהֶל מוֹעֵד לִפְנֵי הַפָּרֹכֶת׃}
}
{
	\eP{
		And he set the goֹldeֶn altar in the Teֶnt oְf Meֵetiִng in
		front of the pavilion,
	}
}
{
}

\Verse{27}
{
	\hP{וַיַּקְטֵר עָלָיו קְטֹרֶת סַמִּים כַּאֲשֶׁר צִוָּה יְהֹוָה אֶת מֹשֶׁה׃}
}
{
	\eP{
		and he burned the iִnceֶnse of fragrances on it, as YHWH had
		commanded Moses.
	}
}
{
}

\Verse{28}
{
	\hP{וַיָּשֶׂם אֶת מָסַךְ הַפֶּתַח לַמִּשְׁכָּן׃}
}
{
	\eP{And he set the cover of the entrance of the Tabernacle.}
}
{
}

\Verse{29}
{
	\hP{וְאֵת מִזְבַּח הָעֹלָה שָׂם פֶּתַח מִשְׁכַּן אֹהֶל מוֹעֵד וַיַּעַל עָלָיו אֶת הָעֹלָה וְאֶת הַמִּנְחָה כַּאֲשֶׁר צִוָּה יְהֹוָה אֶת מֹשֶׁה׃}
}
{
	\eP{
		And he set the altar of burnt offering at the entrance of
		Tabernacle of the Teֶnt oְf Meֵetiִng, and he offered up the
		burnt offering and the grain offering on it, as YHWH had
		commanded Moses.
	}
}
{
}

\Verse{30}
{
	\hP{וַיָּשֶׂם אֶת הַכִּיֹּר בֵּין אֹהֶל מוֹעֵד וּבֵין הַמִּזְבֵּחַ וַיִּתֵּן שָׁמָּה מַיִם לְרחְצָה׃}
}
{
	\eP{
		And he set the basin between the Teֶnt oְf Meֵetiִng and the
		altar, and he put water for washing there.
	}
}
{
}

\Verse{31}
{
	\hP{וְרָחֲצוּ מִמֶּנּוּ מֹשֶׁה וְאַהֲרֹן וּבָנָיו אֶת יְדֵיהֶם וְאֶת רַגְלֵיהֶם׃}
}
{
	\eP{
		And Moses and Aַaroֹn and his sons washed their hands and
		their feet from it.
	}
}
{
}

\Verse{32}
{
	\hP{בְּבֹאָם אֶל אֹהֶל מוֹעֵד וּבְקרְבָתָם אֶל הַמִּזְבֵּחַ יִרְחָצוּ כַּאֲשֶׁר צִוָּה יְהֹוָה אֶת מֹשֶׁה׃}
}
{
	\eP{
		When they came to the Teֶnt oְf Meֵetiִng and when they came
		forward to the altar they would wash, as YHWH had
		commanded Moses.
	}
}
{
}

\Verse{33}
{
	\hP{וַיָּקֶם אֶת הֶחָצֵר סָבִיב לַמִּשְׁכָּן וְלַמִּזְבֵּחַ וַיִּתֵּן אֶת מָסַךְ שַׁעַר הֶחָצֵר וַיְכַל מֹשֶׁה אֶת הַמְּלָאכָה׃}
}
{
	\eP{
		And he set up the courtyard all around the Tabernacle and
		the altar, and he put on the cover of the courtyard’s
		gate. And Moses finished the work.
	}
}
{
}

\Verse{34}
{
	\hP{וַיְכַס הֶעָנָן אֶת אֹהֶל מוֹעֵד וּכְבוֹד יְהֹוָה מָלֵא אֶת הַמִּשְׁכָּן׃}
}
{
	\eP{
		And the cloud covered the Teֶnt oְf Meֵetiִng, and YHWH’s
		glory filled the Tabernacle.
	}
}
{
}

\Verse{35}
{
	\hP{וְלֹא יָכֹל מֹשֶׁה לָבוֹא אֶל אֹהֶל מוֹעֵד כִּי שָׁכַן עָלָיו הֶעָנָן וּכְבוֹד יְהֹוָה מָלֵא אֶת הַמִּשְׁכָּן׃}
}
{
	\eP{
		And Moses was not able to come into the Teֶnt oְf Meֵetiִng,
		because the cloud had settled on it and YHWH’s glory
		filled the Tabernacle.
	}
}
{
}

\Verse{36}
{
	\hP{וּבְהֵעָלוֹת הֶעָנָן מֵעַל הַמִּשְׁכָּן יִסְעוּ בְּנֵי יִשְׂרָאֵל בְּכֹל מַסְעֵיהֶם׃}
}
{
	\eP{
		And when the cloud was lifted from on the Tabernacle, the
		children of Israel would travel—in all their travels—
	}
}
{
}

\Verse{37}
{
	\hP{וְאִם לֹא יֵעָלֶה הֶעָנָן וְלֹא יִסְעוּ עַד יוֹם הֵעָלֹתוֹ׃}
}
{
	\eP{
		and if the cloud would not be lifted, then they would not
		travel until the day that it would be lifted.
	}
}
{
}

\Verse{38}
{
	\hP{כִּי עֲנַן יְהֹוָה עַל הַמִּשְׁכָּן יוֹמָם וְאֵשׁ תִּהְיֶה לַיְלָה בּוֹ לְעֵינֵי כל בֵּית יִשְׂרָאֵל בְּכל מַסְעֵיהֶם׃}
}
{
	\eP{
		Because YHWH’s cloud was on the Tabernacle by day, and
		fire would be in it at night, before the eyes of all the
		house of Israel in all their travels.
    }
}
{
}

%		COMMENTARY ON THE TORAH
%		After the rush of forces in Exodus, Leviticus is a
%		rather tranquil book. Following the mighty acts of God in
%		Exodus, Leviticus is disproportionately the words of God.
%		There is very little narrative—barely three chapters out
%		of twenty-seven in the book—and no poetry (with the
%		possible exception of a few isolated verses). More than
%		the other four of the Five Books of Moses, Leviticus
%		pictures the deity speaking. Consistent with this is the
%		absence of movement in Leviticus. Unlike Genesis and
%		Exodus, the entire book of Leviticus is set in one place:
%		at the foot of Mount Sinai. In Exodus, the universe—the
%		land, waters, and sky—is in disarray. Following that,
%		Leviticus is concerned with orderliness, arrangement.
%		Leviticus pictures the beginning of the solidification—the
%		formation of an identity—of a people through law. The law
%		is significant as a subject of study in itself and in
%		context of the study of history as well. But, also in
%		narrative terms, it is important to see the place of the
%		law in the development of the story of the people: laws of
%		how to behave toward other human beings, laws that bring
%		identity (what to eat, religious ceremonies, holidays).
%		The laws are not merely listed. They are narrated as words
%		of YHWH to Moses (and Aaron) and thus are presented as an
%		integrated component of the account. In Leviticus the
%		knowledge of God, which is an explicit, continuing
%		component in Exodus, likewise resides in the background.
%		The term “to know” (Hebrew yd‘) never occurs in Leviticus
%		in the sense of knowledge of God. YHWH now is known to
%		Israel. This is assumed. Impressively, the story stops
%		being about Moses in Leviticus. The laws are conveyed
%		through him, certainly, but the book is about the laws
%		themselves and the fact of their being revealed to the
%		people. The only stories are primarily about Aaron and his
%		sons. More broadly, Leviticus does not develop
%		personalities. It is about institutions, not individuals.
%		And this, too, contributes to the reduction of tension in
%		Leviticus following the great confrontations of Exodus.
%		All of this contributes to the change of feeling of this
%		book. Leviticus is both integral to the context of the
%		books that precede and follow it and, at the same time,
%		anomalous in character. Leviticus Unlike in Exodus, there
%		are few miracles portrayed in Leviticus, yet the book,
%		rich in background like the two books that precede it, is
%		pervaded with the presence of the divine. The continuing
%		scene has already been set in Exodus: A column of cloud or
%		fire is visible at all times. A Tent of Meeting, the
%		Tabernacle, stands outside the camp. Moses goes in and out
%		of the tent, his face veiled and then unveiled, bringing
%		laws from God. The book begins with God speaking from the
%		Tabernacle, thus establishing the Tabernacle from the
%		outset as the visible, present entity among the people
%		through which the divine word enters the world. The
%		cosmic/miraculous thus looms, and a channel between it and
%		the earthly operates.
